\documentclass[aspectratio=169]{beamer}
\usetheme{Madrid}
\usecolortheme{default}
\usepackage[utf8]{inputenc}
\usepackage{amsmath, amssymb}
\usepackage{graphicx}
\usepackage{booktabs}
\usepackage{tcolorbox}

\definecolor{UCBblue}{RGB}{0, 51, 153}
\setbeamercolor{palette primary}{bg=UCBblue,fg=white}
\setbeamercolor{palette secondary}{bg=UCBblue,fg=white}
\setbeamercolor{palette tertiary}{bg=UCBblue,fg=white}
\setbeamercolor{titlelike}{bg=UCBblue,fg=white}
\setbeamercolor{item}{fg=UCBblue}

\title[Nivelación - Sesión 1]{Taller de Nivelación / Introducción a la Mecatrónica:\\ Álgebra Lineal Aplicada a la Teoría de Redes}
\subtitle{Módulos 1 \& 2: Fundamentos LVK y Sistematización Matricial}
\author[M.Sc. Ing. B. Quiroga Turdera]{M.Sc. Ing. Bernardo Quiroga Turdera}
\institute[UCB Tarija]{Carrera de Ingeniería Mecatrónica\\ Universidad Católica Boliviana "San Pablo"}
\date{28 de Julio de 2026}

\begin{document}

\begin{frame}
    \titlepage
\end{frame}

\begin{frame}{Agenda de la Sesión}
    \tableofcontents
\end{frame}

\section{El Puente Físico-Matemático}

\begin{frame}{1. Ley de Voltajes de Kirchhoff (LVK) y Matrices}
    \textbf{Ley de Voltajes de Kirchhoff ($\sum V = 0$):} Es la manifestación macroscópica del principio de conservación de la energía. En un lazo cerrado (malla), la energía inyectada por las fuentes (elevación de potencial) debe igualar a la energía disipada por las resistencias (caída de potencial).
    
    \vspace{0.3cm}
    \begin{itemize}
        \item \textbf{La Matriz ($\mathbf{R}$):} El \textit{mapa topológico} del hardware. Almacena la oposición geométrica al flujo de electrones.
        \item \textbf{El Vector de Estado ($\mathbf{I}$):} La incógnita de ingeniería; el flujo de carga real (Amperios) que energizará los actuadores.
        \item \textbf{El Sistema ($\mathbf{R}\mathbf{I} = \mathbf{V}$):} Representa el \textit{punto de equilibrio energético} determinista de la red.
    \end{itemize}
\end{frame}

\section{Formulación Algebraica (LVK Paso a Paso)}

\begin{frame}{2. Algoritmo Analítico de Lazo Cerrado}
    Antes de ensamblar tensores matriciales, debemos deducir la ecuación algebraica de la malla ($\sum V_{\text{caídas}} = \sum V_{\text{fuentes}}$):
    
    \vspace{0.2cm}
    \begin{enumerate}
        \item \textbf{Asignación de Estado:} Trazar una corriente de malla en sentido horario ($\circlearrowright$).
        \item \textbf{Caídas de Tensión (Ley de Ohm):}
        \begin{itemize}
            \item Resistencia \textit{propia} (solo pasa su corriente $I_1$): $+ R \cdot I_1$
            \item Resistencia \textit{compartida} con otra malla ($I_2$): $+ R \cdot (I_1 - I_2)$
        \end{itemize}
        \item \textbf{Fuentes Activas (Balance):} Observar por qué terminal ingresa la flecha de la corriente a la batería:
        \begin{itemize}
            \item Entra por negativo (trazo corto): Se anota \textbf{Positivo}.
            \item Entra por positivo (trazo largo): Se anota \textbf{Negativo}.
        \end{itemize}
    \end{enumerate}
\end{frame}

\section{Sistematización Algorítmica}

\begin{frame}{3. Ensamblaje de la Matriz de Resistencias $\mathbf{R}$}
    Una vez dominada la formulación algebraica, reordenamos los términos para construir el modelo matricial (Paradigma CDIO: Concebir $\rightarrow$ Diseñar).
    
    \vspace{0.3cm}
    \begin{tcolorbox}[colback=white,colframe=UCBblue,title=Reglas Topológicas Directas]
    \begin{itemize}
        \item \textbf{Impedancia Propia ($R_{kk}$):} Suma de todas las resistencias del lazo $k$. Se ubica en la diagonal principal y es \textbf{estrictamente positiva} ($R_{kk} > 0$).
        \item \textbf{Impedancia Mutua ($R_{kj}$):} Resistencia compartida en la frontera topológica entre el lazo $k$ y el $j$. Se ubica fuera de la diagonal y va \textbf{precedida por un signo negativo}.
        \item \textbf{Firma Estructural:} La matriz resultante es \textbf{simétrica} ($R_{kj} = R_{jk}$).
    \end{itemize}
    \end{tcolorbox}
\end{frame}

\section{Determinantes y Limitaciones}

\begin{frame}{4. El Determinante y la Regla de Sarrus}
    El determinante ($\det(\mathbf{R})$) es un escalar físico que valida si nuestro circuito tiene un único punto de equilibrio.
    
    \vspace{0.3cm}
    \textbf{Expansión Visual a $3\times3$ (Regla de Sarrus):}
    \begin{equation*}
        \det(\mathbf{R}) = (R_{11}R_{22}R_{33} + \dots) - (R_{13}R_{22}R_{31} + \dots)
    \end{equation*}
    
    \begin{alertblock}{Limitación Crítica en Ingeniería}
        \textbf{La Regla de Sarrus es una heurística visual exclusiva para matrices $3\times3$.} Carece de validez matemática para dimensiones $4\times4$ o superiores. Por ello, delegaremos el cálculo al entorno computacional (Python/NumPy).
    \end{alertblock}
\end{frame}

\section{Singularidad y Transición Computacional}

\begin{frame}{5. Diagnóstico de Falla: La Matriz Singular}
    Si $\det(\mathbf{R}) = 0$, el sistema matemático colapsa. Físicamente, esto significa:
    \begin{itemize}
        \item \textbf{Fuentes en Conflicto:} Baterías en paralelo forzando distinto potencial.
        \item \textbf{Cortocircuito Ideal:} Lazo cerrado sin resistencia disipativa ($I \to \infty$).
    \end{itemize}
\end{frame}

\begin{frame}{6. Transición al Entorno Computacional}
    \textbf{Fase CDIO: Implementar y Operar}
    \begin{enumerate}
        \item Inicialicen el entorno Jupyter Notebook (\texttt{Notebook\_01} y \texttt{Notebook\_02}).
        \item Transcriban el modelo matricial a un tensor (\texttt{np.array}).
        \item Validaremos las propiedades de la red:
        \begin{itemize}
            \item Simetría Estructural: \texttt{assert np.array\_equal(R, R.T)}
        \end{itemize}
        \item Sistematizaremos el determinante con \texttt{np.linalg.det(R)} e implementaremos un manejador de excepciones (\texttt{try-except LinAlgError}) para aislar singularidades (cortocircuitos).
    \end{enumerate}
    
    \vspace{0.3cm}
    \begin{center}
        \textit{Fin de la fundamentación analítica. Iniciamos el despliegue computacional.}
    \end{center}
\end{frame}

\end{document}
